%! ~~~ Packages Setup ~~~ 
\documentclass[]{../../../../tex/classes/styledArticle}
\usepackage{../../../../tex/packages/hebrewSupport}
\usepackage{../../../../tex/packages/mathShortcuts}
\usepackage{../../../../tex/packages/theoremsSupport}
\usepackage{../../../../tex/packages/enfancysections}

%! ~~~ Document ~~~

\author{\en{Shahar Perets}}
\title{\textit{לינארית 2א 21}}
\date{24 ביוני 2025}
\begin{document}
	\maketitle
	\theo{אם $A \in M_n(\F)$ נורמלית, אזי קיים פולינום $f(x) \in \R[x]$ כך ש־$A^* = f(A)$}
	
	דיברנו על המשפט הזה בשבוע שעבר, אבל ההוכחה הייתה קצת עקומה אז נחזור עליה. 
	\begin{proof}
		עבור $A$ נורמלית מהמשפט הספקטרלי קיים בסיס אורתונורמלי מלכסן ולכן קיימת $P$ הפיכה כך ש־$P\op APש = \diag(\lg_1 \dots \lg_n)$. לכן $P\op A^* P = \diag(\bar \lg_1 \dots \bar \lg_n)$. נשתמש במשפט לפיו יש פולינום $f \in \R[x]$ כך ש־$f(x_i) = \bar x_i$ ובפרט בעבור $x_i = \lg_i$ קיים פולינום עבורו $f(\lg_i) = \bar \lg_i$. אזי 
		\[ f(\diag(\lg_1 \dots \lg_n)) = \diag(f(\lg_1) \dots f(\lg_n)) = \diag(\bar \lg_1 \dots \bar \lg_n) \]
		\[ f(P\op A P) = P\op f(A) P = P \op A^* P \implies f(A) = A^* \]
		עוד נבחין ש־$\deg f = n - 1$. 
	\end{proof}
	
	ננסה להבין מי הן $A \in M_2(\R)$ שהן נורמליות. מעל $\C$ הן פשוט לכסינות. נבחין ש־: 
	\begin{align*}
		A = \pms{a & b \\ c & d}, \ A^* = f(A) \in \R_{\le1}[x] = \ag A + BI, \ A = \ag A^T + \bg I = \ag(\ag A + \bg I) + \bg I \\ \implies \pms{a & b \\ c & d} = \pms{\ag^2 a + \bg(\ag + 1) & \ag^2 b \\ \ag^2 c & \ag^2 d + \bg (\ag + 1)} \\
		\begin{cases}
			(b \land c \neq 0 ) \quad \ag = 1 \implies  A = A + 2 \bg I \implies \bg = 0, \ A = A^T \cancel{+ \bg I} \\
			(b \land c \neq 0) \quad \ag = -1 \implies A^T = - A  +\bg I  \implies A + A^T = \pms{\bg & 0 \\ 0 & \bg} \implies b = -c, \ A = \pms{a & b \\ -b & a} \\
			(b \lor c = 0) \implies A = \pms{a & 0 \\ 0 & d}
		\end{cases}
	\end{align*}
	המקרה השני – זה פשוט סיבובים, אבל בניפוח (כי הדטרמיננטה היא $a^2 - b^2$). 
	
	בכל מקרה, מסקנה מהמשפט הקודם. 
	\theo{אם $T \co V \to V$ נורמלית, אז $\exists f \in \R[x] \co f(T) = T^*$. }
	\begin{proof}
		נבחר בסיס א''נ $A^* = [T^*]_B, \impliedby A = [T]_B$. כבר הוכחנו שאם $T$ נורמלית אז $A$ נורמלית ולכן מהמשפט הקודם קיים $f$ מתאים כך ש־$[T^*]_B = A^* = f(A) = f([T]_B) = [f(T)]_B$. סה''כ $[T^*]_B = [f(T)]_B$ ומחח''ע העברת בסיס $T^* = f(T)$ כדרוש. 
	\end{proof}
	
	אם $T \co V \to V$ ט''ל, $U, W \subseteq V$ תמ''וים $T$־איוונריאנטי כך ש־$U \oplus W = B$. אם $\bc$ בסיס של $V$, כאשר קישא של הבסיס הוא הבסיס של $U$ אז: 
	\[ [T]_\bc = \pms{[T|_U]_\bc &  \\  & [T|_W]_\bc} \]
	בפרט בעבור ניצבים $U \subseteq V \implies V = U \oplus U^{\perp}$. ניעזר בכך כדי להוכיח את המשפט הבא: 
	\theo{אם $U \subseteq V$ תמ''ו אינוו' ביחס ל־$T$ אז $U^{\perp}$ הוא $T^*$־אינוו'. }\begin{proof}
		יהי $w \in U^{\perp}$. רוצים להראות $T^*w \in U^{\perp}$. יהי $u \in U$ אז: 
		\[ \mut{T^* w}{u} = \mut{w}{Tu} = \mut{w}{u'}, \ u' \in U \implies \mut{w}{u'} = 0 \quad \top \]
	\end{proof}
	\theo{בעבור $T \co V \to V$ נורמלית, אם היא $U$־אינו' אז גם $T^*$ הוא $U$־אינו'}
	\begin{proof}
		נבחין ש־$T^* = f(T)$ כלשהו, וכן $U$ הוא $T$־איוו' ולכן $U$ הוא $f(T)$־איוו' וכאן די גמרנו את ההוכחה. 
	\end{proof}
	מסימטריות $U^{\perp}$ הוא $T^*$, מהמשפט גם $(T^*)^*$ איונ' ולכן $T$־אינו'. 
	
	\theo{יהי $V$ מעל $\R$ מ''ו וכן $T \co V \to V$ ט''ל. אז קיים $U \subseteq V$ שהוא $T$־איונ' וממדו לכל היותר $T$. }
	\textit{הערה: }מעל $\C$ ``זה מטופש'' כי הפולינום מתפרק (ואז המרחב העצמי יקיים את זה). 
	\begin{proof}
		נפרק ל־$m_T(x)$ מינימלי ו־$g(x)$ גורם אי־פריק כך ש־$m_T(x) = g(x)h(x)$. לכל $g$ אי פריק ב־$\R$ הוא לינארי  הוא ממעלה $2$, מהמשפט היסודי של האגלברה ומהעובדה ש־$m_T(x) = 0 \implies m_T(\bar x) = 0$. 
		\begin{itemize}
			\item אם $g$ לינארי אז יש ע''ע ממשי של $T$ מה שנותא $U$ (שנפרש ע''י הו''ע) ממד $1$. 
			\item אם $\deg g = 2$ כמובן שניתן להניח $g$ מתוקן. ז''א $g(x) = x^2 + ax + b$ אז $g(T)$ אינו הפיך (מלמת החלוקה לפולינום מינימלי) כלומר $\exists 0 \neq v \in \ker g(T)$. 
			לכן: 
			\[ (T^2 + aT + bI)v = 0 \implies T^2v = -a T v - bv \]
			ולכן $U = \Sp(v, Tv)$ תמ''ו עם ממד לכל היותר 2 וגם נשמר תחת $T$. 
		\end{itemize}
	\end{proof}
	
	\textit{הערה: }בעבור נורמלית הטענה נכונה ללא תלות במשפט היסודי של האלגברה. 
	
	לכן, בעבור נורמליות, מהמשפט הספקטרלי ומטענות קודמות, עבור $T \co V \to V$ ממשית קיים בסיס א''נ $\bc$ של $V$ שבעבורו המטריצה המייצגת של $T$ היא מטריצת בלוקים $2 \times 2$ מצורה של $\pms{a & b \\ -b & a}$: 
	\[ [T]_\bc = \diag\cl{\pms{a_1 & b_1 \\ -b_1 & a_1} \cdots \pms{a_k & b_k \\ -a_k & b_k}, \ \lg_1 \cdots \lg_m} \]
	כאשר כמובן $2k + m = n$. 
	
	\section{\en{Unitary matrixes}}
	\defi{יהי $V$ ממ''פ. אז $T \co V \to V$ תקרא \textit{אוניטרית} (אם $\F = \C$) או \textit{אורתוגונלית} (אם $\F = \R$) אם $T^*T = I$ או במילים אחרות $T^* = T\op$. }
	ברור שט''ל כזו היא נורמלית. 
	\textbf{דוגמה. }עבור $T_\tg$ הסיבוב ב־$\tg$ מעלות, במישור $\R^2$, אז $(T_\tg)^* = T_{-\tg} = T\op_{\tg}$.  
	\textbf{דוגמה. }עבור $T$ שיקוף מתקים $T^2 = I$ וכן $T^* = T$ וסה''כ $T^* = T = T\op$. 
	
	\theo{התנאים הבאים על $T \co V \to V$ שקולים: 
	\begin{itemize}
		\item \hfil $T^* = T\op$
		\item \hfil $\forall v, u \co \mut{Tv}{Tu} = \mut{v}{u}$
		\item $T$ מעבירה כל בסיס א''נ של $V$ לבסיס א''נ של $V$
		\item $T$ מעבירה בסיס א''נ אחד של $V$ לבסיס א''נ של $V$ (כלומר, מספיק להראות שקיים בסיס יחידה שעובר לבסיס אחר). 
		\item \hfil $\forall v \in V \co \norm{Tv} = \norm{v}$
	\end{itemize}}
	כלומר: היא משמרת זווית (העתקה פנימית) וגודל. במילים אחרות, היא משמרת העתקה פנימית. 
	\begin{proof}נפרק לרצף גרירות
		\begin{enumerate}
			\item[$1 \to 2$] \hfil $T^* = T\op \implies \mut{Tv}{Tu} = \mut{v}{T^*Tu} = \mut{v}{u}$
			\item[$2 \to 3$] נאמר ש־$(v_1 \dots v_n)$ א''נ. צ.ל. $(Tv_i)_{i = 1}^{n}$ א''נ. לשם כך נצטרך להוכיח את שני התנאים – החלק של האורתו והחלק של הנורמלי. בשביל שניהם מספיק להוכיח ש־: 
				$\mut{Tv_i}{Tv_j} = \mut{v_i}{v_j}  = \dg_{ij}$
			\item[$3 \to 4$]טרוויאלי
			\item[$4 \to 5$]יהי $(v_1 \dots v_n)$ בסיס א''נ כך ש־$(Tv_1 \dots Tv_n)$ א''נ. אז: 
			\begin{align*}
				v = \sumni \ag_i v_i \implies &\norm{v}^2 = \mut{\sumni \ag_i v_i}{\sumni \ag_i v_i} = \sumni |\ag_i|^2 \\
				&\norm{Tv}^2 = \mut{T\cl{\sumni \ag_i v_i}}{T\cl{\sumni \ag_i v_i}} = \mut{\sumni \ag_i T(v_i)}{\sumni \ag_i T(v_i)} = \sum |\ag_i|^2
			\end{align*}
			\item[$5 \to 1$]מניחים $\forall v \in V \co \norm{Tv} = \norm{v}$. ידועות השקילויות הבאות: 
			\[ T^* = T\op \iff T^*T = I \iff T^*T - I = 0 \]
			בעבר ראינו את הטענה הבאה: נניח ש־$S$ צמודה לעצמה וכן ש־$\forall v \co \mut{Sv}{v} = 0$, אז $S = 0$. במקרה הזה: 
			\[ S := T^*T - I \implies S^* = (T^*T - I)^* = (T^*T)^* - I^* = T^*(T^*)^* - I =S \implies S^* = S \]
			והיא אכן צמודה לעצמה. עוד נבחין ש־: 
			\[ \mut{Sv}{v} = \mut{(T^* T - I)v}{v} =\mut{T^*Tv}{v} - \mut{v}{v} = \mut{Tv}{Tv} - \mut{v}{v} = \norm{Tv}^2 - \norm{v}^2 = 0 \]
			השוויון האחרון נכון מההנחה היחידה שלנו ש־$\norm{Tv} = \norm{v}$. סה''כ $TT^* - I = 0$. סה''כ הוכחנו $TT^* - I = 0$ שזה שקול ל־$T^* = T\op$ מהשקילויות לעיל כדרוש. 
		\end{enumerate}
	\end{proof}
	
	\theo{תהי $T \co V \to V$ אוני'אורתו', ו־$\lg$ ע''ע של $T$. אז $|\lg| = 1$} \begin{proof}
		יהי $v$ ו''ע של הע''ע $\lg$. אז: 
		\[ |\lg|^2 \mut{v}{v} = \lg\bar\lg \mut{v}{v} = \mut{Tv}{Tv} = \mut{v}{v} \]
	\end{proof}
	\defi{תהי $A \in M_n(\F)$. אז $A$ תיקרא \textit{אוניטרית}/\textit{אורתוגונלית} אם $A^* = A\op$. }
	\theo{אוניטרית: $A\ol{A^T} = I$}
	\theo{אורתוגונלית: $AA^T = I$}
	\textit{הערה: }אוניטרית בה מלשון unit – היא שומרת על הגודל, על וקטורי היחידה (ה־unit vectors). 
	\theo{יהי $\bc$ בסיס א''נ של $V$ ו־$T \co V \to V$ אז $T$ אוניטרית/אורתוגונלית אמ''מ $A = [T]_B$ אוניטרית/אורתוגונלית. }
	\begin{proof}
		\[ AA^* = [T]_\bc[T^*]_\bc = [TT^*]_\bc, I = AA^* \iff [TT^*]_\bc = I \iff TT^* = I \]
	\end{proof}
	
	הנושאים האחרונים למבחן – צורה קאנונית של העתקה נורמלית מעל הממשיים, וכן העתקות אורתוגונליות ואוניטריות. ללא צורה קאנונית של העתקה אורתוגונלית. לכן ההרצאות בשיעורים הבאים לא נכנסות לחומר. תודה איראן. 
	
	\textit{הערה. }\textit{איזומטריה} היא העתקה שמשמרת גדלים, ואיזומטריה אורתוגונלית היא פשוט אוניטרית. משום מה זה שם שמדברים עליו בע''פ אבל לא הגדירו מסודר. 
	\ndoc
\end{document}
